\scanpage{190}
% manuscript position: top
\mspage{106}
% manuscript position: 1
s'insérant ds le diagramme commutatif
% manuscript position: 2–7
\diagram{{\scriptsize\begin{tikzcd}[ampersand replacement=\&,column sep=0.8em,row sep=large]
\underline{\mathrm{Hom}}_{!\pi}(C,\unsure{\mathcal E|E}) \arrow[r,"\psi","{(\approx)}"'] \arrow[d,hook,"\text{inc. pl. fid.}"] \& \underline{\mathrm{Hom}}_{!}(\operatorname{Ind}_\pi(C),\mathcal E)\arrow[r]\arrow[d,dashed,"\text{pl. fid}"] \& \underline{\mathrm{Hom}}_{!\pi}(\operatorname{Ind}_\pi(\unsure{C}),\mathcal E)\\
\underline{\mathrm{Hom}}(C,\mathcal E)\arrow[r,hook,"{\substack{L\\\approx}}","{\substack{\deleted{\text{pl. fid}}\\\deleted{\text{\unsure{d'image ess. formée|prolonge un foncteur}}}\\\deleted{\text{des foncteurs qui}}\\\deleted{\text{commute aux}}\\\deleted{\text{petites }\varinjlim\text{ filtrantes}}}}"'] \& \underline{\mathrm{Hom}}_{!\,\mathrm{filt}}(\operatorname{Ind}(\deleted{\unsure{C}}C),\mathcal E)\arrow[ur,"{R\ (\text{restriction})}"'] \& {}
\end{tikzcd}}}
% manuscript position: left/3
(4.4.16.\unsure{6|5})
% manuscript position: 8
[cf SGA 4 I 8.7.3]

% manuscript position: 9–10
4.5. \uline{Filtration cardinale canonique d'une} \inserted{catégorie} \inserted{$\mathcal U$-}\uline{paratopos}.

% manuscript position: 11–14
\uline{Définition} 4.5.1. On appelle \uline{\inserted{catégorie-$\mathcal U$-}para\deleted{\illegible{1}}topos} une catégorie \deleted{qui est} \inserted{$M$} \inserted{$\mathcal U$-}pseudo\deleted{\illegible{1}}topos et \deleted{\illegible{1}} accessible, i.e qui satisfait aux conditions équivalentes du théorème 4.3.\unsure{5|6} (p 70).
% manuscript position: 15–25
On dit \deleted{la description} \unsure{que|qu'}, c'est une \deleted{$\pi$-}catégorie $\pi$-\inserted{$\mathcal U$-}para\deleted{\illegible{1}}topos si elle est \inserted{\deleted{\unsure{pleine}}}$\pi$-accessible, i.e. si elle admet une petite sous-catégorie génératrice dont les objets sont $\pi$-\unsure{accessible|accessibles}, ou ce qui revient au même, si $M$ est équivalente à une catégorie de la forme $\operatorname{Ind}'_\pi(M_\pi)$, où $M_\pi$ est une \inserted{petite} catégorie [\inserted{ou : une cat.} équivalente à une petite catégorie], et stable par \deleted{petites} $\varinjlim$ de cardinal $\leq\pi$. (Et on peut prendre \deleted{\unsure{pour}} $M_\pi$ la sous-catégorie pleine de $M$ formée des objets
% manuscript position: foot
% archive sheet number: 189
